\question{1.7}{}
Let $N$ denote the \textbf{north pole} $(0,\cdots,0,1)\in \mathbb{S}^n\subseteq \mathbb{R}^{n+1}$, and let $S$ denote the \textbf{south pole} $(0,\cdots,0,-1)$. Define the \textbf{stereographic projection} $\sigma:\mathbb{S}^n\setminus \{N\}\to \mathbb{R}^n$ by
\begin{align}
    \sigma(x_1,\cdots,x_{n+1})=\frac{(x_1,\cdots,x_n)}{1-x_{n+1}}.
\end{align}
Let $\tilde{\sigma}(x)=-\sigma(-x)$ for $x\in \mathbb{S}^n\setminus \{S\}$. 
\begin{itemize}
    \item [(a)] For any $x\in \mathbb{S}^n\setminus \{N\}$, show that $\sigma(x)=u$, where $(u,0)$ is the point where the line through $N$ and $x$ intersects the linear subspace where $x_{n+1}=0$. Similarly, show that $\tilde{\sigma}(x)$ is the point where the line through $S$ and $x$ intersects the same subspace. (For this reason, $\tilde{\sigma}$ is also called $\textbf{stereographic projection from the south pole}$.)  
    \item [(b)] Show that $\sigma$ is bijective and 
    \begin{align}
        \sigma^{-1}(u_1,\cdots,u_n)=\frac{(2u_1,\cdots,2u_n, |u|^2-1)}{|u|^2+1}.
    \end{align}
    \item [(c)] Compute the transition map $\tilde{\sigma}\circ \sigma^{-1}$ and verify that the atlas consisting of the two charts $(\mathbb{S}^n\setminus \{N\},\sigma)$ and $(\mathbb{S}^n\setminus \{S\},\tilde{\sigma})$ defines a smooth structure on $\mathbb{S}^n$. (The coordinates defined by $\sigma$ and $\tilde{\sigma}$ are called \textbf{stereographic coordinates}.)
    \item [(d)] Show that this smooth structure is the same as the one defined in Example~1.31.
\end{itemize}

\textbf{Example 1.31} We put a smooth structure on $\mathbb{S}^n$ as follows. For each $i=1,\cdots,n+1$, let $(U_i^\pm,\phi_i^\pm)$ denote the graph coordinate charts, where 
\begin{align}
    U_i^\pm=\{(x_1,\cdots,x_{n+1})\in \mathbb{S}^n: \pm x_i>0\},\quad \phi_i^\pm(x_1,\cdots,x_{n+1})=(x_1,\cdots,\hat{x}_i,\cdots,x_{n+1}),
\end{align}
where the hat indicates that the $i$-th coordinate is omitted. 


For any distinct indices $i,j$, the transition map $\phi_i^\pm\circ (\phi_j^\pm)^{-1}$ is easily computed. In the case $i<j$, we get 
\begin{align}
    \phi_i^\pm\circ (\phi_j^\pm)^{-1}(u_1,\cdots,u_{n})=(u_1,\cdots,\hat{u}^{i},\cdots,\pm\sqrt{1-|u|^2},\cdots,u_{n})
\end{align}
with the $\pm\sqrt{1-|u|^2}$ in the $j$-th position, and a similar formula holds for $i>j$. In the case $i=j$, an even simpler formula holds:
\begin{align}
    \phi_i^\pm\circ (\phi_i^\mp)^{-1}(u_1,\cdots,u_{n})=\mathrm{ID}_{\mathbb{B}^n}.
\end{align}
Thus, the collection of charts $\{(U_i^\pm,\phi_i^\pm)\}$ is a smooth atlas, and so defines a smooth structure on $\mathbb{S}^n$.

\answer{}
\begin{itemize}
    \item [(a)]
\end{itemize}
Consider the line 
\begin{align}
    \ell(t)=N+t(x-N)=t(x_1,\cdots,x_n,x_{n+1}-1)+(0,\cdots,0,1)=(tx_1,\cdots,tx_n,t(x_{n+1}-1)+1).
\end{align}
Obviously, $\ell(0)=N$ and $\ell(1)=x$. We want to find the point where this line intersects the hyperplane $x_{n+1}=0$. Setting the last coordinate of $\ell(t)$ equal to zero, we have
\begin{align}
    t(x_{n+1}-1)+1=0\implies t=\frac{1}{1-x_{n+1}}.
\end{align}
Plugging this value of $t$ into the first $n$ coordinates of $\ell(t)$, we get
\begin{align}
    \ell\left(t=\frac{1}{1-x_{n+1}}\right)=\frac{(x_1,\cdots,x_n)}{1-x_{n+1}}.
\end{align}
This is exactly the definition of $\sigma(x)$. A similar argument shows that $\tilde{\sigma}(x)$ is the point where the line through $S$ and $x$ intersects the hyperplane $x_{n+1}=0$. More explicitly, consider the line
\begin{align}
    \tilde{\ell}(t)=S+t(x-S)=t(x_1,\cdots,x_n,x_{n+1}+1)+(0,\cdots,0,-1)=(tx_1,\cdots,tx_n,t(x_{n+1}+1)-1).
\end{align}
Setting the last coordinate of $\tilde{\ell}(t)$ equal to zero, we have
\begin{align}
    t(x_{n+1}+1)-1=0\implies t=\frac{1}{1+x_{n+1}}.
\end{align}
Plugging this value of $t$ into the first $n$ coordinates of $\tilde{\ell}(t)$, we get
\begin{align}
    \tilde{\ell}\left(t=\frac{1}{1+x_{n+1}}\right)=\frac{(x_1,\cdots,x_n)}{1+x_{n+1}}.
\end{align}
This is exactly the definition of $\tilde{\sigma}(x)$.

\begin{itemize}
    \item [(b)] 
\end{itemize}
In order to show that $\sigma$ is bijective, we will show that it is both injective and surjective. For injectivity, suppose that $\sigma(x)=\sigma(y)$ for some $x,y\in \mathbb{S}^n\setminus \{N\}$. Then we have 
\begin{align}
    \frac{(x_1,\cdots,x_n)}{1-x_{n+1}}=\frac{(y_1,\cdots,y_n)}{1-y_{n+1}}.
\end{align}
Hence, we have 
\begin{align}
    x_i= \frac{1-x_{n+1}}{1-y_{n+1}}y_i\quad \text{for all } i=1,\cdots,n.
\end{align}
Consider the last coordinate of $x$ and $y$. Since both $x$ and $y$ are on the unit sphere, we have 
\begin{align}
    1-x_{n+1}^2=\sum_{i=1}^{n}x_i^2=\sum_{i=1}^{n}\left(\frac{1-x_{n+1}}{1-y_{n+1}}y_i\right)^2=\left(\frac{1-x_{n+1}}{1-y_{n+1}}\right)^2\sum_{i=1}^{n}y_i^2=\left(\frac{1-x_{n+1}}{1-y_{n+1}}\right)^2(1-y_{n+1}^2).
\end{align}
This implies that
\begin{align}
    &1-x_{n+1}^2=\left(\frac{1-x_{n+1}}{1-y_{n+1}}\right)^2(1-y_{n+1}^2)\\
    \implies &(1-x_{n+1})(1+x_{n+1})=\frac{(1-x_{n+1})^2}{(1-y_{n+1})^2}(1-y_{n+1})(1+y_{n+1})\\
    \implies &(1+x_{n+1})(1-y_{n+1})=(1-x_{n+1})(1+y_{n+1})\\
    \implies &1+x_{n+1}-y_{n+1}-x_{n+1}y_{n+1}=1-x_{n+1}+y_{n+1}-x_{n+1}y_{n+1}\\
    \implies &2x_{n+1}=2y_{n+1}\implies x_{n+1}=y_{n+1}.
\end{align}
If $x_{n+1}=y_{n+1}$, then we have
\begin{align}
    x_i= \frac{1-x_{n+1}}{1-y_{n+1}}y_i=y_i\quad \text{for all } i=1,\cdots,n.
\end{align}
Thus, we have shown that $x=y$, and hence $\sigma$ is injective.

Next, in order to show that $\sigma$ is surjective, let $u\in \mathbb{R}^n$ be arbitrary. We want to find $x\in \mathbb{S}^n\setminus \{N\}$ such that $\sigma(x)=u$. Let $x=(x_1,\cdots,x_n,x_{n+1})$. Then we have
\begin{align}
    \sigma(x)=u\implies \frac{(x_1,\cdots,x_n)}{1-x_{n+1}}=u\implies (x_1,\cdots,x_n)=(1-x_{n+1})u=(u^1(1-x_{n+1}),\cdots,u^n(1-x_{n+1})).
\end{align}
With the constraint that $x\in \mathbb{S}^n$, we have
\begin{align}
    1=x_1^2+\cdots+x_n^2+x_{n+1}^2=(1-x_{n+1})^2|u|^2+x_{n+1}^2\implies 1=(1-x_{n+1})^2|u|^2+x_{n+1}^2.
\end{align}
This is a quadratic equation in $x_{n+1}$, which we can simplify as follows:
\begin{align}
    0=(1-x_{n+1})^2|u|^2+x_{n+1}^2-1=(|u|^2+1)x_{n+1}^2-2|u|^2x_{n+1}+|u|^2-1.
\end{align}
Solving for $x_{n+1}$, we have
\begin{align}
    x_{n+1}=\frac{2|u|^2\pm \sqrt{4|u|^4-4(|u|^2+1)(|u|^2-1)}}{2(|u|^2+1)}=\frac{|u|^2\pm 1}{|u|^2+1}.
\end{align}
Since we want $x\in \mathbb{S}^n\setminus \{N\}$, we must have $x_{n+1}\neq 1$, which means that we must take the negative sign in the above expression. Hence, we have 
\begin{align}
    x_{n+1}=\frac{|u|^2-1}{|u|^2+1}.
\end{align}
Now that we have $x_{n+1}$, we can find $x_i$ for $i=1,\cdots,n$:
\begin{align}
    x_i=u^i(1-x_{n+1})=u^i\left(1-\frac{|u|^2-1}{|u|^2+1}\right)=u^i\left(\frac{2}{|u|^2+1}\right)=\frac{2u^i}{|u|^2+1}.
\end{align}
Thus, we have found $x\in \mathbb{S}^n\setminus \{N\}$ such that $\sigma(x)=u$, and hence $\sigma$ is surjective. Therefore, $\sigma$ is bijective. Beside, we have shown that
\begin{align}
    \sigma^{-1}(u_1,\cdots,u_n)=\frac{(2u_1,\cdots,2u_n, |u|^2-1)}{|u|^2+1}.
\end{align}

\begin{itemize}
    \item [(c)]
\end{itemize}
Consider the transition map $\tilde{\sigma}\circ \sigma^{-1}$. For any $u\in \sigma(\mathbb{S}^n\setminus \{N\})$, we have
\begin{align}
    \tilde{\sigma}\circ \sigma^{-1}(u)&=\tilde{\sigma}\left(\frac{(2u_1,\cdots,2u_n, |u|^2-1)}{|u|^2+1}\right)
    =\frac{\frac{1}{|u|^2+1}(2u_1,\cdots,2u_n)}{1+\frac{|u|^2-1}{|u|^2+1}}
    =\frac{(2u_1,\cdots,2u_n)}{2 |u|^2}\\
    &=\frac{(u_1,\cdots,u_n)}{|u|^2}.
\end{align}
This is a smooth map, and hence the atlas consisting of the two charts $(\mathbb{S}^n\setminus \{N\},\sigma)$ and $(\mathbb{S}^n\setminus \{S\},\tilde{\sigma})$ defines a smooth structure on $\mathbb{S}^n$. Beside, the transition map $\sigma\circ \tilde{\sigma}^{-1}$ is also smooth, since it is the inverse of a smooth map. Therefore, the atlas consisting of the two charts $(\mathbb{S}^n\setminus \{N\},\sigma)$ and $(\mathbb{S}^n\setminus \{S\},\tilde{\sigma})$ defines a smooth structure on $\mathbb{S}^n$.

\begin{itemize}
    \item [(d)]
\end{itemize}
To show that the smooth structure defined by stereographic coordinates is the same as the one defined in Example~1.31, we need to show that the transition maps between the two atlases are smooth.
Let $(U_i^\pm,\phi_i^\pm)$ be a chart from the atlas defined in Example~1.31, and let $(\mathbb{S}^n\setminus \{N\},\sigma)$ be the chart defined by stereographic coordinates. We need to show that the transition map $\phi_i^\pm\circ \sigma^{-1}$ is smooth. For any $u\in \sigma(U_i^\pm \cap \mathbb{S}^n\setminus \{N\})$, we have 
\begin{align}
    \phi_i^\pm\circ \sigma^{-1}(u)&=\phi_i^\pm\left(\frac{(2u_1,\cdots,2u_n, |u|^2-1)}{|u|^2+1}\right)\\
    &=\left(\frac{2u_1}{|u|^2+1},\cdots,\frac{2u_{i-1}}{|u|^2+1}, \hat{\frac{2u_i}{|u|^2+1}}, \frac{2u_{i+1}}{|u|^2+1}, \cdots, \frac{2u_n}{|u|^2+1},\frac{|u|^2-1}{|u|^2+1}\right).
\end{align}
for $i=1,\cdots,n$. For $i=n+1$, we have
\begin{align}
    \phi_{n+1}^\pm\circ \sigma^{-1}(u)&=\phi_{n+1}^\pm\left(\frac{(2u_1,\cdots,2u_n, |u|^2-1)}{|u|^2+1}\right)\\
    &=\left(\frac{2u_1}{|u|^2+1},\cdots,\frac{2u_n}{|u|^2+1}\right).
\end{align}
In both cases, the transition map is smooth. Next, we need to show that the transition map $\sigma\circ (\phi_i^\pm)^{-1}$ is smooth. For any $v\in \phi_i^\pm(U_i^\pm \cap \mathbb{S}^n\setminus \{N\})$, we have
\begin{align}
    \sigma\circ (\phi_i^\pm)^{-1}(v)&=\sigma\left((v_1,\cdots,v_{i-1},\underbrace{\pm\sqrt{1-|v|^2}}_{\text{the $i$-th coordinate}},v_i,\cdots,v_n)\right)\\
    &=\frac{1}{1-v_n}\left(v_1,\cdots,v_{i-1},\pm\sqrt{1-|v|^2},v_i,\cdots,v_{n-1}\right),
\end{align}
for $i=1,\cdots,n$. For $i=n+1$, we have
\begin{align}
    \sigma\circ (\phi_{n+1}^\pm)^{-1}(v)&=\sigma\left((v_1,\cdots,v_n,\underbrace{\pm\sqrt{1-|v|^2}}_{\text{the $(n+1)$-th coordinate}})\right)\\
    &=\frac{1}{1-\pm\sqrt{1-|v|^2}}\left(v_1,\cdots,v_n\right)\\
    &=\frac{1}{1\mp\sqrt{1-|v|^2}}\left(v_1,\cdots,v_n\right).
\end{align}
Both of these transition maps are smooth. Therefore, the smooth structure defined by stereographic coordinates is the same as the one defined in Example~1.31.





\clearpage
\question{1.9}{}
\textbf{Complex projective $n$-space}, denoted by $\mathbb{CP}^n$, is the set of all 1-dimensional complex-linear subspaces of $\mathbb{C}^{n+1}$, with the quotient topology inherited
from the natural projection $\pi:\mathbb{C}^{n+1}\setminus \{0\}\to \mathbb{CP}^n$. Show that $\mathbb{CP}^n$ is a
compact $2n$-dimensional topological manifold, and show how to give it a smooth structure analogous to the one we constructed for $\mathbb{RP}^n$. (We use the correspondence
\begin{align}
    (x^1+iy^1,\cdots,x^{n+1}+iy^{n+1}) \longleftrightarrow (x^1,\cdots,x^{n+1},y^1,\cdots,y^{n+1})
\end{align}
to identify $\mathbb{C}^{n+1}$ with $\mathbb{R}^{2n+2}$.)

\answer{}
